
\exo 

De 1985 à 2010, les températures moyennes relevées au mois de janvier dans la
ville de Luxembourg sont les suivantes.

-4,4 ; 1,1 ; -4,3 ; 4,3 ; 2,1 ; 2,4 ; 0,8 ; 0,6 ; 2,7 ; 2,4 ; 1,0 ; -1,2 ;
-2,4 ; 2,1 ; 2,9 ; 1,7 ; 1,8 ; 1,6 ; 0,3 ; 1,5 ; 3,2 ; 0,9 ; 6,1 ; 5,1 ;
-0,7 ; -0,9.
\textit{\footnotesize(source : http://www.statistiques.public.lu)}
\begin{enumerate}
    \item Déterminer les quartiles et la médiane de cette série.
    \item  Interpréter les valeurs trouvées en écrivant deux phrases sans
        utiliser les mots:
        \begin{multicols}{2}
        \begin{enumerate}
        	\item quartile; \item médiane.
        \end{enumerate}
        \end{multicols}
\end{enumerate}



\exo Dépense durant les soldes

Le montant des
dépenses (en euros) de chaque client lors d'une journée de soldes  a été relevé et
trié  dans le tableau ci-dessous où les fréquences sont exprimées en pourcentage. 
\vspace{0.5em}

\begin{center}
\begin{tabular}{c|*{6}{ >{\centering} m{1.5cm}}}
\hline
  Classe &    [10;30[ &   [30;50[ & [50;70[ &    [70;90[ & [90;110[&[110;130]\tabularnewline 
\hline
Fréquences &   15 & 25 & 10 &   20 & 10 & 20\tabularnewline 
\hline
\end{tabular} 
\end{center}


\begin{enumerate}
\item Construire le polygone des fréquences cumulées croissantes.
\item 
Déterminer par lecture graphique, une approximation
\begin{multicols}{3}
\begin{itemize}
 \item de la médiane;
 \item du premier quartile;
\item du troisième quartile.
\end{itemize}
\end{multicols}
\item Interpréter ces résultats.
\item 
Déterminer une approximation de la moyenne. \\
La lecture graphique est-elle possible? \\
Interpréter ce résultat.
\end{enumerate}


\exo Les cigognes sont passées

Les tailles sont exprimées en centimètre.

à la maternité \og  Beaux jours\fg
Sur la totalité du mois de janvier 2012, il y a eu \\57 nouveau-nés à la
maternité
\og  Beaux jours\fg.\\ Leur taille est donnée dans le tableau ci-dessous.
\smallskip

\begin{tabular}{c*{14}{>{\centering}m{0.9cm}}}\hline
 Taille & 46&47,5&48&48,5&49&49,5&50& 50,5&51&51,5&52&52,5&53\tabularnewline \hline
 Effectifs & 1&2&3&5&5&7&9&8&7&5&2&2&1\tabularnewline \hline
\end{tabular}

 \begin{enumerate}
        \item Calculer la moyenne puis la médiane des tailles de ces 57~nouveau-nés en précisant la démarche.
        \item Calculer le pourcentage de nouveau-nés ayant une taille 
inférieure ou égale à 49 cm.

Donner la réponse arrondie à 0,1 près.

\item Parmi toutes ces tailles, déterminer la plus petite taille $t$ 
telle qu'au moins les trois quarts des nouveau-nés aient une taille inférieure 
ou égale à $t$~cm. \\
Quel paramètre de la série des tailles a été ainsi
trouvé?
\end{enumerate}



